\documentclass[a4paper,12pt]{article}
\usepackage[T2A]{fontenc}
\usepackage[utf8]{inputenc}
\usepackage[russian]{babel}
\usepackage{amsmath,amssymb,mathtools}
\usepackage{helvet}
\renewcommand{\familydefault}{\sfdefault}
\usepackage{icomma}
\usepackage{geometry}
\geometry{left=20mm,right=20mm,top=20mm,bottom=19mm}
\usepackage{microtype}
\usepackage{tikz}
\usetikzlibrary{arrows.meta,positioning,calc,shapes.geometric}
\usepackage{pgfplots}
\pgfplotsset{compat=1.18}
\usepackage{tabularx,booktabs}
\usepackage{enumitem}
\usepackage{hyperref}
\usepackage{parskip}
\usepackage{fancyhdr}
\usepackage{xcolor}
\definecolor{accent}{HTML}{14777A}
\definecolor{darkaccent}{HTML}{164E51}
\definecolor{muted}{HTML}{687675}
\definecolor{gridcolor}{HTML}{E3E9E8}
\hypersetup{hidelinks,pdfauthor={Лёвин Артём Александрович},pdftitle={Чек-лист: метод опорных точек для прямой, параболы и гиперболы}}
\setlength{\parindent}{0pt}
\setlength{\parskip}{6pt}
\setlist[itemize]{leftmargin=6mm,itemsep=2pt,topsep=3pt}
\setlist[enumerate]{leftmargin=8mm,itemsep=3pt,topsep=3pt}
\renewcommand{\arraystretch}{1.25}
\pagestyle{fancy}
\fancyhf{}
\renewcommand{\headrulewidth}{0pt}
\fancyfoot[L]{\footnotesize\textcolor{muted}{А.\,А. Лёвин \quad / \quad Анна Баннова}}
\fancyfoot[R]{\footnotesize\textcolor{muted}{\thepage}}
\newcommand{\checkitem}{\(\square\)\enspace}
\newcommand{\sectionline}{\par\vspace{-4pt}\noindent\textcolor{accent}{\rule{\linewidth}{0.55pt}}\par\vspace{3pt}}
\newcommand{\accentbold}[1]{\textcolor{darkaccent}{\textbf{#1}}}
\begin{document}
\thispagestyle{empty}
\begin{center}
\vspace*{23mm}
{\large\textcolor{accent}{\textbf{МАТЕМАТИКА / ПРОФИЛЬНЫЙ ЕГЭ}}}\par\vspace{22mm}
{\fontsize{33}{39}\selectfont\bfseries ЧЕК-ЛИСТ}\par\vspace{11mm}
{\fontsize{22}{28}\selectfont\bfseries Метод опорных точек}\par\vspace{5mm}
{\large Прямая \(\cdot\) парабола \(\cdot\) гипербола}\par
\vspace{19mm}
\textcolor{accent}{\rule{0.58\linewidth}{0.7pt}}\par\vspace{16mm}
{\large\bfseries Анна Баннова}\par\vspace{4mm}
{\normalsize 25.10.26}\par\vspace{11mm}
{\normalsize\textbf{Лёвин Артём Александрович}\\[2mm]
эксклюзивно для Анны Банновой}
\end{center}
\vfill
\begin{center}
\begin{tikzpicture}
\draw[accent!70,line width=1.15pt]
 (0,0) .. controls (2.7,1.05) and (5.2,-.9) .. (7.9,.1)
       .. controls (10.8,.95) and (13,-.7) .. (15.8,.25);
\end{tikzpicture}
\end{center}
\clearpage

\section*{01. Основа метода}
\sectionline
\checkitem Различать координаты точки \(A(x_0;y_0)\): \(x_0\)~--- аргумент, \(y_0=f(x_0)\)~--- значение функции.\par
\checkitem Выбирать две удобные точки графика с \emph{различными} абсциссами.\par
\checkitem Подставлять обе пары координат в формулу и получать систему двух уравнений.\par
\checkitem Находить неизвестные коэффициенты методом подстановки.\par
\checkitem Восстанавливать формулу \(f(x)\) и только затем отвечать на вопрос задачи.\par

\subsection*{Почему подходят разные графики}
Пусть известен \emph{вид} функции, но неизвестны два коэффициента. Для прямой, параболы с известным старшим коэффициентом и гиперболы действует один принцип:
\[
\boxed{\ (x_1;y_1),\ (x_2;y_2)\ \longrightarrow\
\begin{cases}y_1=f(x_1),\\y_2=f(x_2)\end{cases}
\ \longrightarrow\ \text{коэффициенты}\ \longrightarrow\ f(x)\ }.
\]
Удобнее брать точки с \(x=0,\pm1\), если такие точки принадлежат графику. \textbf{Исключение:} у функции \(k/x+a\) аргумент \(x=0\) недопустим.

\section*{02. Прямая: \(f(x)=kx+b\)}
\sectionline
На графике прямой отмечены \(A(-1;-3)\) и \(B(3;4)\). Требуется найти \(f(-5)\).
\begin{center}
\begin{tikzpicture}
\begin{axis}[
 width=11.5cm,height=5.7cm,axis lines=middle,xmin=-3,xmax=4.2,
 ymin=-6,ymax=6.5,xtick={-3,-2,-1,0,1,2,3,4},ytick={-6,-4,-2,0,2,4,6},
 grid=both,grid style={gridcolor,very thin},minor tick num=0,
 xlabel={$x$},ylabel={$y$},tick label style={font=\footnotesize},clip=false]
\addplot[accent,very thick,domain=-2.6:4,samples=2]{1.75*x-1.25};
\addplot[only marks,mark=*,mark size=2.7pt,accent] coordinates {(-1,-3) (3,4)};
\node[below right,font=\small] at (axis cs:-1,-3) {$A(-1;-3)$};
\node[above left,font=\small] at (axis cs:3,4) {$B(3;4)$};
\end{axis}
\end{tikzpicture}
\end{center}
\textbf{1. Подстановка координат:}
\[
\begin{cases}
-3=k\cdot(-1)+b,\\
4=k\cdot3+b
\end{cases}
\ \Longleftrightarrow\
\begin{cases}-3=-k+b,\\4=3k+b.\end{cases}
\]
\textbf{2. Решение системы:} из первого уравнения \(b=k-3\). Тогда
\[
4=3k+(k-3)\ \Longleftrightarrow\ 4k=7\ \Longleftrightarrow\ k=\frac74,
\qquad b=\frac74-3=-\frac54.
\]
\textbf{3. Формула и ответ на вопрос задачи:}
\[
f(x)=\frac74x-\frac54,\qquad
f(-5)=\frac74\cdot(-5)-\frac54=-\frac{40}{4}=\boxed{-10}.
\]
\clearpage

\section*{03. Парабола: \(f(x)=2x^2+bx+c\)}
\sectionline
У функции известен коэффициент при \(x^2\). Нужно восстановить \(b\) и \(c\) по точкам \(A(0;-4)\), \(B(1;1)\).
\begin{center}
\begin{tikzpicture}
\begin{axis}[
 width=11.6cm,height=7.1cm,axis lines=middle,xmin=-3.1,xmax=2.3,
 ymin=-5.8,ymax=5.6,xtick={-3,-2,-1,0,1,2},ytick={-5,-4,-3,-2,-1,0,1,2,3,4,5},
 grid=both,grid style={gridcolor,very thin},minor tick num=0,
 xlabel={$x$},ylabel={$y$},tick label style={font=\footnotesize},clip=true]
\addplot[accent,very thick,domain=-2.8:2.0,samples=230]{2*x^2+3*x-4};
\addplot[only marks,mark=*,mark size=2.8pt,accent] coordinates {(0,-4) (1,1)};
\node[below right,font=\small] at (axis cs:0,-4) {$A(0;-4)$};
\node[above left,font=\small] at (axis cs:1,1) {$B(1;1)$};
\end{axis}
\end{tikzpicture}
\end{center}
\textbf{1. Составляем систему:}
\[
\begin{cases}
-4=2\cdot0^2+b\cdot0+c,\\
1=2\cdot1^2+b\cdot1+c.
\end{cases}
\]
\textbf{2. Находим коэффициенты:}
\[
c=-4,\qquad 1=2+b-4\ \Longrightarrow\ b=3.
\]
\textbf{3. Восстанавливаем функцию:}
\[
\boxed{f(x)=2x^2+3x-4}.
\]
\accentbold{Выбор точек.} Точка с \(x=0\) сразу даёт \(c=f(0)\); точка с \(x=1\) существенно упрощает второе уравнение. Направление ветвей влияет на знак старшего коэффициента, но сам метод остаётся прежним.

\vspace{4mm}
\textbf{Самопроверка:} \(f(0)=-4\), \(f(1)=2+3-4=1\). Обе точки согласованы с найденной формулой.
\clearpage

\section*{04. Гипербола: \(f(x)=\dfrac{k}{x}+a\)}
\sectionline
По графику известны \(A(-3;0)\) и \(B(1;4)\). Найдём \(k\) и \(a\).
\begin{center}
\begin{tikzpicture}
\begin{axis}[
 width=11.6cm,height=6.7cm,axis lines=middle,xmin=-5.4,xmax=4.2,
 ymin=-6.1,ymax=6,xtick={-5,-4,-3,-2,-1,0,1,2,3,4},
 ytick={-6,-4,-2,0,2,4,6},grid=both,grid style={gridcolor,very thin},minor tick num=0,
 xlabel={$x$},ylabel={$y$},tick label style={font=\footnotesize},clip=true]
\addplot[accent,very thick,domain=-5.3:-0.43,samples=230]{3/x+1};
\addplot[accent,very thick,domain=0.52:4.1,samples=230]{3/x+1};
\addplot[only marks,mark=*,mark size=2.8pt,accent] coordinates {(-3,0) (1,4)};
\node[below left,font=\small] at (axis cs:-3,0) {$A(-3;0)$};
\node[above right,font=\small] at (axis cs:1,4) {$B(1;4)$};
\end{axis}
\end{tikzpicture}
\end{center}
\textbf{1. Подставляем точки.} Учтём \(x\ne0\):
\[
\begin{cases}0=\dfrac{k}{-3}+a,\\[1mm]4=\dfrac{k}{1}+a\end{cases}
\Longleftrightarrow
\begin{cases}0=-\dfrac{k}{3}+a,\\[1mm]4=k+a.\end{cases}
\]
\textbf{2. Выразим \(a=4-k\) и подставим:}
\[
0=-\frac{k}{3}+4-k.
\]
Чтобы убрать дробь, умножаем \emph{каждое слагаемое} на \(3\):
\[
0=-k+12-3k\ \Longleftrightarrow\ 4k=12
\ \Longleftrightarrow\ k=3,\qquad a=4-3=1.
\]
\textbf{3. Искомая формула:}
\[
\boxed{f(x)=\frac3x+1,\quad x\ne0}.
\]
\clearpage

\section*{05. Алгебраические приёмы и типичные ошибки}
\sectionline
\subsection*{«Странствующий минус» у дроби}
Знак минус можно перенести из знаменателя перед дробью или в числитель:
\[
\frac{u}{-v}=-\frac{u}{v}=\frac{-u}{v},\qquad v\ne0.
\]
Если числитель представляет собой сумму, минус относится \emph{ко всему числителю}:
\[
\frac{x+1}{-5}=-\frac{x+1}{5}=\frac{-(x+1)}5
=\frac{-x-1}{5}.
\]
\textbf{Ошибка:} \(\displaystyle\frac{x+1}{-5}=\frac{-x+1}{5}\)~--- неверно: перед единицей тоже должен появиться минус.

\subsection*{Как избавляться от дробей в системе}
При умножении \emph{обеих частей} уравнения на число нужно умножить на него \emph{каждое слагаемое}:
\[
-4=-\frac{k}{2}+a
\quad\xRightarrow{\ \cdot2\ }\quad -8=-k+2a.
\]
На занятии при подстановке разбиралась также пара точек \(A(-2;-4)\), \(B(2;-2)\) для \(f(x)=k/x+a\):
\[
\begin{cases}-4=-\dfrac{k}{2}+a,\\[1mm]-2=\dfrac{k}{2}+a\end{cases}
\ \Longleftrightarrow\
\begin{cases}-8=-k+2a,\\-4=k+2a.\end{cases}
\]
Из второго уравнения \(k=-4-2a\). \textbf{Подставляем со скобками:}
\[
-8=-(-4-2a)+2a=4+4a
\ \Longrightarrow\ a=-3,\quad k=2.
\]
Таким образом, \(\displaystyle f(x)=\frac2x-3\). Проверка: \(f(-2)=-4\), \(f(2)=-2\).

\subsection*{Что чаще всего мешает найти коэффициенты}
\begin{itemize}
\item \textbf{Переставили координаты:} в точке \((x;y)\) сначала идёт \(x\), затем \(y\).
\item \textbf{Потеряли знак:} при раскрытии \(-(-4-2a)\) получается \(4+2a\).
\item \textbf{Подставили без скобок:} вместо коэффициента всегда берите найденное выражение в скобки.
\item \textbf{Домножили только дробь:} число перед скобкой или справа от знака равенства тоже нужно умножить.
\item \textbf{Забыли исходный вопрос:} после нахождения коэффициентов восстановите \(f(x)\), затем найдите требуемое значение функции.
\item \textbf{Взяли точку гиперболы при \(x=0\):} аргумент \(x=0\) запрещён для \(k/x+a\).
\end{itemize}
\clearpage

\section*{06. Алгоритм метода опорных точек}
\sectionline
\vspace{6mm}
\begin{center}
\begin{tikzpicture}[
 node distance=10mm,>=Latex,
 term/.style={ellipse,draw=darkaccent,line width=.85pt,fill=accent!6,align=center,text width=8.6cm,minimum height=10mm,font=\small},
 processnode/.style={rectangle,rounded corners=2pt,draw=accent,line width=.85pt,align=center,text width=10.2cm,minimum height=13mm,font=\small,inner sep=4pt},
 ask/.style={diamond,aspect=2.75,draw=darkaccent,line width=.85pt,fill=accent!6,align=center,text width=4.4cm,inner sep=1pt,font=\small},
 side/.style={rectangle,rounded corners=2pt,draw=accent,line width=.85pt,align=center,text width=3.3cm,minimum height=12mm,font=\small,inner sep=3pt},
 arr/.style={-Latex,line width=.85pt,draw=accent}]
\node[term] (a) {Есть график и вид формулы с двумя неизвестными коэффициентами};
\node[processnode,below=of a] (b) {Выбрать две удобные точки \((x_1;y_1)\), \((x_2;y_2)\)};
\node[processnode,below=of b] (c) {Составить систему \(y_1=f(x_1)\), \(y_2=f(x_2)\)};
\node[ask,below=9mm of c] (d) {Есть дроби в системе?};
\node[side,right=9mm of d] (e) {Домножить уравнения, чтобы убрать дроби};
\node[processnode,below=20mm of d] (f) {Выразить один коэффициент через другой};
\node[processnode,below=of f] (g) {Подставить выражение \textbf{в скобках}; найти оба коэффициента};
\node[processnode,below=of g] (h) {Записать \(f(x)\), проверить по двум исходным точкам};
\node[term,below=of h] (i) {Выполнить действие, указанное в условии, и записать ответ};
\draw[arr] (a)--(b);\draw[arr] (b)--(c);\draw[arr] (c)--(d);
\draw[arr] (d)-- node[left,font=\footnotesize]{нет} (f);
\draw[arr] (d)-- node[above,font=\footnotesize]{да} (e);
\draw[arr] (e.south) |- ([yshift=6mm]f.north) -- (f.north);
\draw[arr] (f)--(g);\draw[arr] (g)--(h);\draw[arr] (h)--(i);
\end{tikzpicture}
\end{center}
\vfill
\noindent\textcolor{muted}{При работе с гиперболой учтите ограничение \(x\ne0\). Проверка двух точек в завершении метода выявляет большинство ошибок в знаках и коэффициентах.}
\end{document}
